\textbf{范围：\CurrentSegmentRange。} 目标是建立清晰的三拍律动，让上行音型逐步积累，再准确收住。

\SegmentScore

\subsubsection*{第一轮：三个小单元}
\begin{enumerate}
\item \textbf{弱起至第 4 小节：认音与路线。} 先按谱面指法练上行音型，找准临时升降号和换位位置。跨谱表的音仍按谱面声部分工处理，不能仅凭音符所在谱表判断用哪只手。
\item \textbf{第 5--7 小节：保持均匀。} 先不加渐强，把每个八分音符弹匀，再加入谱面的渐强。手臂随音区移动，不用突然加速制造推进感。
\item \textbf{第 8--9 小节：高点与停顿。} 单独找准高音落点，数足休止；第 9 小节之后的安静也属于乐句。
\end{enumerate}

\subsubsection*{第二轮：分手、合手、连段}
先分手确认音高、节奏和指法，再慢速合手，检查和弦与旋律的对齐。初学时先少用或不用踏板，听清音符后再按谱面加入。
可尝试以四分音符每分钟 60--72 拍起步；若仍不稳，就继续放慢。连续三遍无停顿、无明显错音，再小幅提速。
最后将第 8--9 小节与下一段第 10--11 小节连接，练习从引子转入轻巧主题的变化。
